\section{결론 및 한계}
\label{sec:limitations}

본 연구의 목적은 단순히 다음 탈락 팀을 맞추는 데 있지 않다. 핵심은 PUBG
텔레메트리를 팀 이동, 안전 구역 수축, 교전 이벤트, 주변 팀 관계를 포함하는
페이즈 수준 위험 상태로 표현했을 때, 그 표현이 실제 생존 과정의 신호를
보존하는지 검증하는 것이다. 21,824개 경기에서 페이즈 정렬 STDB 표현은 전략
점수 요약보다 더 많은 위험 정보를 담았고, 동적 팀 그래프는 특징 기반 모델
위에 작지만 통계적으로 신뢰할 수 있는 관계 신호를 추가하였다.

따라서 이 논문의 주된 의미는 새로운 대형 예측 모델을 제안했다는 데보다,
경기 후 집계 통계로 소비되던 PUBG 텔레메트리를 배틀로얄의 공간 수축과 팀 간
상호작용을 포함한 시공간 생존 과정으로 재정식화했다는 데 있다. 표현
ablation에서는 이 재정식화의 이득이 가장 크게 나타났고, graph ablation은
관계 구조가 그 위에 추가 신호를 제공함을 보였다.

한계도 분명하다. 첫째, 본 연구는 예측 연구이지 인과 연구가 아니다. 관측된
근접, 교전, 공동 압력 관계가 다음 페이즈 탈락 위험 예측에 도움이 된다는
결과가, 특정 관계를 조정하면 생존이 인과적으로 변한다는 뜻은 아니다. 둘째,
텔레메트리는 관측 가능한 위치, 상태, 이벤트만 포함한다. 지형, 건물
내부/외부, 가시선, 음성 의사소통, 플레이어 숙련도와 의도처럼 생존에 중요한
요인은 직접 관측하지 못한다. 셋째, 팀 중심점과 페이즈 그래프는 복잡한 squad
움직임을 압축한다. 팀 분산과 최근접 적 특징으로 일부 보완하지만, 같은 팀
중심점이라도 대형이나 엄폐 상태가 다를 수 있다. 넷째, 경기 단위 무작위
분할은 같은 데이터 분포 안의 일반화를 평가하므로, 패치, 시즌, 지역, 대회
환경이 바뀌었을 때의 일반화까지 보장하지는 않는다. 후반 페이즈에는 남은 팀과
위험 셀이 적어 추정 변동성과 생존자 편향도 커진다. 마지막으로 그래프 구성은
고정된 거리 및 압력 임계값을 사용하므로, 향후 연구에서는 임계값 민감도와
분포 변화 평가를 함께 수행해야 한다.
